% Part: set-theory
% Chapter: z
% Section: nat

\documentclass[../../../include/open-logic-section]{subfiles}

\begin{document}

\olfileid{sth}{z}{nat}
\olsection{আমাদের স্বাভাবিক সংখ্যা নির্বাচন}

%ধরি, অনানুষ্ঠানিক নীতি \stagesinf{}-কে ন্যায্যতা দেওয়া যায়।
%তবু ওই নীতি থেকে পাওয়া নির্দিষ্ট স্বতঃসিদ্ধটি নিয়েও
%কয়েকটি কথা বলা দরকার।

\olref[infinity-again]{defnomega}-এ আমরা ``স্বাভাবিক
সংখ্যা'' কথাটির অর্থ স্পষ্টভাবে নির্ধারণ করেছি। এই
নির্ধারণ কীভাবে বুঝবে? এটি কোনো অধিবিদ্যাগত দাবি নয়;
কেবল নির্দিষ্ট কয়েকটি সেটকে স্বাভাবিক সংখ্যা
\emph{হিসেবে ধরার} সিদ্ধান্ত। এভাবে ভাবার কারণগুলি
\olref[sfr][rel][ref]{sec}, \olref[sfr][arith][ref]{sec} এবং
\olref[sfr][infinite][dedekindsproof]{sec}-এ কিছুটা দেখেছি।
এখন সেগুলি আরও স্পষ্ট করা যায়।

আমাদের অসীমতার স্বতঃসিদ্ধটি \citet{VonNeumann1925}-এর
পথ অনুসরণ করে। তবে চাইলে এর বদলে নিচের স্বতঃসিদ্ধটিও
নিতে পারতাম:

\begin{defish}
\emph{জার্মেলোর \citeyear{Zermelo1908Untersuchungen}-এর
অসীমতার স্বতঃসিদ্ধ।} এমন একটি সেট $A$ আছে, যাতে
$\emptyset \in A$ এবং $(\forall x \in A)\{x\} \in A$।
\end{defish}

আমাদের (ফন নয়মান-অনুপ্রাণিত) অসীমতার স্বতঃসিদ্ধের
পরিবর্তে জার্মেলোরটি নিলেও একটি ডেডেকিন্ড-অসীম সেট,
আর তাই একটি ডেডেকিন্ড বীজগঠন, সমানভাবেই পাওয়া যেত।
জার্মেলোর পদ্ধতিতে আমাদের বীজগঠনের বিশিষ্ট সদস্যও
হতো $\emptyset$ ($0$-এর বিকল্প); কিন্তু এক-এক চিত্রণটি
$x \mapsto x \cup \{x\}$-এর বদলে $x \mapsto \{x\}$
হতো। এর সবচেয়ে সহজ ফল হলো, জার্মেলো $2$ হিসেবে
$\{\{\emptyset\}\}$ ধরেন, আর আমরা (ফন নয়মানের মতো)
$2$ হিসেবে $\{\emptyset, \{\emptyset\}\}$ ধরি।

একটি অসীমতার স্বতঃসিদ্ধ ছেড়ে অন্যটি নেব কেন? প্রধান
ব্যবহারিক কারণ, ফন নয়মানের পদ্ধতি ট্রান্সফাইনাইট সংখ্যা
নিয়েও বেশ ভালোভাবে কাজ করে।
\olref[ordinals][]{chap} থেকে আমরা তা দেখব। কিন্তু
শুধু \emph{পাটিগণিত করার} দিক থেকে দুটি পদ্ধতিই
সমান উপযোগী। তাই কেউ যদি বলে যে স্বাভাবিক সংখ্যাগুলি
\emph{আসলে} সেট, সঙ্গে-সঙ্গেই প্রশ্ন আসে:
\emph{কোন সেটগুলি?}

এই নির্দিষ্ট প্রশ্নটি \citet{Benacerraf1965} সুপরিচিত
করে তোলেন। তবে মনে রাখতে হবে, এরকম ব্যাপার আমরা
আগেও বহুবার দেখেছি; এটি কেবল তার সবচেয়ে পরিচিত
দৃষ্টান্ত। মূল কথাটি হলো: সেটতত্ত্বের সাহায্যে আমরা
বিভিন্ন ``স্বজ্ঞাগত'' ধরনের সত্তাকে
\emph{প্রতিনিধিত্ব} করতে পারি—বাস্তব, মূলদ, পূর্ণ ও স্বাভাবিক
সংখ্যা যেমন, তেমনি ক্রমযুগল, অপেক্ষক ও সম্পর্ককেও।
কিন্তু সেটতত্ত্ব কখনো প্রতিনিধিত্বের
\emph{একমাত্র} উপায় নির্ধারণ করে দেয় না। সব ক্ষেত্রেই
বিকল্প থাকে, যা অনায়াসেই একই কাজ করতে পারত।

\end{document}
